% Section 2.2, from lectures/L02/appendix/c_arduino_via_microchip_studio.md.

\section{Arduino via Microchip Studio}
\label{c2:sec:studio}\label{c2:app:c}

\subsection{Why this workflow}
\label{c2:c:1}
If you don't have a dedicated ISP programmer, you can still flash a Microchip Studio build onto the
board: the Arduino bootloader already sits on the chip, and the Arduino IDE already knows the exact
\code{avrdude} command line needed to talk to it over your board's COM port. The trick is simply to
borrow that command line and point it at Microchip Studio's \code{.hex} output instead of the
Arduino IDE's own. Everything below builds toward one reusable ``External Tool'' entry in Microchip
Studio that does exactly that.

\subsection{Wire the circuit}
\label{c2:c:2}
Same layout as the first example of \secref{c2:app:a}: LED anode on D9 (through a
220--470\,\textohm{} resistor to ground, 330\,\textohm{} is a good middle value), button on D13
(one leg to 5\,V, the other to ground through a 10\,k\textohm{} resistor, with D13 connected
between the button and the resistor).

\bookimage{lectures/L02/appendix/images/01_breadboard_wiring.png}{The circuit on the breadboard:
the LED on D9 and the button on D13, with its pull-down resistor.}{c2:fig:wiring}

\noindent Like the diagrams in \secref{c2:app:a}, this one shows an Arduino Uno. If you're using a
Nano instead: same chip, same pin numbering, just a different board outline in the picture.

\subsection{Build and test the sketch in Arduino IDE}
\label{c2:c:3}
This step uses the Arduino libraries (C++, an abstraction layer this book doesn't otherwise use)
purely to get a known-working upload before touching Microchip Studio. Open the Arduino IDE, start
a new sketch (\code{File → New}), and replace its contents with:

\begin{cppcode}
/** GPIO pins. */
#define LED1 9U  // D9.
#define BTN1 13U // D13.

/**
 * @brief Set up system.
 */
void setup()
{
    // Configure LED1 as output.
    pinMode(LED1, OUTPUT);

    // Configure BTN1 as input (the external 10 kOhm resistor pulls it down).
    pinMode(BTN1, INPUT);
}

void loop()
{
    // Read BTN1, enable LED1 if pressed, disable if not.
    if (digitalRead(BTN1)) { digitalWrite(LED1, HIGH); }
    else { digitalWrite(LED1, LOW); }
}
\end{cppcode}

\noindent This lights the LED on D9 while the button on D13 is held. It'll be rebuilt as real C in
\secref{c2:c:7}.

\code{pinMode()}, \code{digitalRead()}, and \code{digitalWrite()} are Arduino's abstraction over
the register-level I/O \secref{c2:app:a} covers directly: \code{pinMode(pin, mode)} sets a pin
\code{INPUT} or \code{OUTPUT} (a \code{DDRx} write), \code{digitalRead(pin)} returns
\code{HIGH}/\code{LOW} for a pin's current level (a \code{PINx} read), and
\code{digitalWrite(pin, value)} sets an output pin \code{HIGH} or \code{LOW} (a \code{PORTx} write).
Underneath, the Arduino core implements all three with exactly the register operations from
\secref{c2:app:a}, just hidden behind a function call and a pin-number lookup table.

Turn on verbose upload logging, since the next step depends on reading \code{avrdude}'s exact
command line out of it: \code{File → Preferences}, then check the box next to ``Show verbose output
during: Upload''.

\bookimage[0.95]{lectures/L02/appendix/images/02_file_menu_preferences.png}{The Arduino IDE's File
menu, with Preferences.}{c2:fig:c02}

\bookimage[0.95]{lectures/L02/appendix/images/03_preferences_verbose_upload.png}{The Preferences
dialog, with verbose output turned on for upload.}{c2:fig:c03}

\subsection{Install the board package, select the board and port}
\label{c2:c:4}
Install the AVR board package if it isn't already: \code{Tools → Board → Boards Manager}, then click
Install on ``Arduino AVR Boards by Arduino'' (skip this if it's already installed, as in the
screenshot below).

\bookimage[0.95]{lectures/L02/appendix/images/04_tools_board_boards_manager_menu.png}{Tools, Board,
Boards Manager.}{c2:fig:c04}

\bookimage[0.95]{lectures/L02/appendix/images/05_boards_manager_installed.png}{The Boards Manager,
with Arduino AVR Boards already installed.}{c2:fig:c05}

\noindent With the board connected, select it: \code{Tools → Board → Arduino AVR Boards → Arduino
Uno}. Then select its port: \code{Tools → Port}, and pick the COM port shown (e.g.\ COM4; the exact
number doesn't matter and can vary by machine).

\bookimage[0.95]{lectures/L02/appendix/images/06_tools_board_arduino_uno.png}{Selecting the Arduino
Uno as the board.}{c2:fig:c06}

\bookimage[0.95]{lectures/L02/appendix/images/07_tools_port_com4.png}{Selecting the board's port,
here COM4.}{c2:fig:c07}

\subsection{Upload and read back avrdude's command line}
\label{c2:c:5}
Click Upload in the top-left corner.

\bookimage[0.95]{lectures/L02/appendix/images/08_upload_button.png}{The Upload button, the arrow in
the top-left corner.}{c2:fig:c08}

\noindent In the Output window (the black pane), scroll up to the line reading ``Compilation
complete.'' and copy the \code{avrdude} command on the line beneath it; it should look something
like:

\begin{console}
"C:\Users\eripih2001\AppData\Local\Arduino15\packages\arduino\tools\avrdude\6.3.0-arduino17/bin/avrdude"
"-CC:\Users\eripih2001\AppData\Local\Arduino15\packages\arduino\tools\avrdude\6.3.0-arduino17/etc/avrdude.conf"
-v -V -patmega328p -carduino "-PCOM4" -b115200 -D
"-Uflash:w:C:\Users\ERIPIH~1\AppData\Local\Temp\arduino-sketch4CE86550DB8430968B53E52972DC68D2/Led1.ino.hex:i"
\end{console}

\bookimage[0.95]{lectures/L02/appendix/images/09_output_window_avrdude_command.png}{The avrdude
command in the Output window, on the line after ``Compilation complete.''}{c2:fig:c09}

\subsection{Turn it into a reusable command}
\label{c2:c:6}
This is the fiddliest part, but it's a one-time setup: reshaping that copied line into a
command/arguments pair Microchip Studio's ``External Tools'' feature can invoke directly.
\begin{enumerate}
  \item Open Notepad, paste the copied path in, and strip out all the quotation marks. Save it as
    \code{Arduino via Microchip Studio.txt}.
  \item After \code{bin/avrdude}, press Enter twice so the rest of the path lands two lines down.
  \item Add the extension \code{.exe} after \code{avrdude} on the top line.
\end{enumerate}

\bookimage[0.95]{lectures/L02/appendix/images/10_notepad_path_pasted.png}{The command in Notepad,
its quotation marks removed and the avrdude path on a line of its own.}{c2:fig:c10}

\bookimage[0.95]{lectures/L02/appendix/images/11_notepad_exe_appended.png}{\code{.exe} added after
\code{avrdude} on the top line.}{c2:fig:c11}

\begin{enumerate}[resume]
  \item Far to the right on the second line, just before \code{-Uflash}, press Enter so the rest
    lands on a new third line.
  \item On the two bottom lines, add quotes around the two file paths (the ones starting with
    \code{C:\\} and ending in \code{.conf} and \code{.hex} respectively). Afterward the two lines
    should look roughly like:
\end{enumerate}

\begin{console}
-C"C:\Users\eripih2001\AppData\Local\Arduino15\packages\arduino\tools\avrdude\6.3.0-arduino17/etc/avrdude.conf" -v -V -patmega328p -carduino -PCOM4 -b115200 -D
-Uflash:w:"C:\Users\ERIPIH~1\AppData\Local\Temp\arduino-sketch4CE86550DB8430968B53E52972DC68D2/Led1.ino.hex":i
\end{console}

\bookimage[0.95]{lectures/L02/appendix/images/12_notepad_arguments_split.png}{The arguments split
before \code{-Uflash}.}{c2:fig:c12}

\bookimage[0.95]{lectures/L02/appendix/images/13_notepad_quotes_added.png}{Quotes added around the
\code{.conf} and \code{.hex} paths.}{c2:fig:c13}

\begin{enumerate}[resume]
  \item Cut the file path on the bottom line (the part after \code{-Uflash}, between the quotes
    ending in \code{.hex}) and paste in \code{$(ProjectDir)Debug\$(TargetName).hex} instead:
    Microchip Studio expands these build-variable placeholders to the actual output path at run
    time, so this line no longer needs to change per project.
  \item Merge the two bottom lines back into one: click at the very start of the last line (right
    before \code{-Uflash}) and press Backspace. Only two lines should remain in total. Save the
    file.
\end{enumerate}

\bookimage[0.95]{lectures/L02/appendix/images/14_notepad_state_before_hex_replace.png}{The
\code{.hex} path replaced by Microchip Studio's build variables.}{c2:fig:c14}

\bookimage[0.95]{lectures/L02/appendix/images/15_notepad_hex_path_replaced.png}{The finished file:
the Command on the top line, the Arguments below it.}{c2:fig:c15}

\noindent You now have two lines: the \code{avrdude.exe} path (the \emph{Command}) and everything
else (the \emph{Arguments}): ready to paste into Microchip Studio's External Tools dialog in
\secref{c2:c:8}.

\subsection{Recreate the program in Microchip Studio}
\label{c2:c:7}
Open Microchip Studio and create a new \textbf{GCC C Executable Project}, named \code{Led1}.

\bookimage[0.95]{lectures/L02/appendix/images/16_new_project_dialog.png}{A new GCC C Executable
Project, named Led1.}{c2:fig:c16}

\noindent In the Device Selection dialog, choose \textbf{ATmega328P}.

\bookimage[0.95]{lectures/L02/appendix/images/17_device_selection.png}{Choosing the ATmega328P in
the Device Selection dialog.}{c2:fig:c17}

\noindent Replace the contents of \file{main.c} with the real-C equivalent of the sketch from
\secref{c2:c:3}, written to this codebase's conventions:

\begin{ccode}
#include <stdint.h>

#include <avr/io.h>

/** GPIO pins. */
#define LED1 1U // D9  -> PORTB1.
#define BTN1 5U // D13 -> PORTB5.

/** GPIO operations. */
#define LED1_ON PORTB |= (1U << LED1)         // Enable LED1.
#define LED1_OFF PORTB &= ~(1U << LED1)       // Disable LED1.
#define BTN1_IS_PRESSED (PINB & (1U << BTN1)) // High if BTN1 is pressed, low otherwise.

/**
 * @brief Set up system.
 */
static void setup(void)
{
    // Configure LED1 as output.
    DDRB = (1U << LED1);

    // Configure BTN1 as input with its pull-up enabled.
    PORTB = (1U << BTN1); // D13 as input, pull-up enabled
}

/**
 * @brief Application entry point.
 *
 * @return 0 on termination of the program (should never occur).
 */
int main(void)
{
    setup();

    while (1)
    {
        // Read BTN1, enable LED1 if pressed, disable if not.
        if (BTN1_IS_PRESSED) { LED1_ON; }
        else { LED1_OFF; }
    }
    return 0;
}
\end{ccode}

\noindent Then build it: \code{Build → Build Solution}.

\bookimage[0.95]{lectures/L02/appendix/images/18_mainc_build_menu.png}{Build, Build
Solution.}{c2:fig:c18}

\noindent The screenshots for this walkthrough were taken against earlier drafts of the program:
of the Arduino sketch in \secref{c2:c:3}--\ref{c2:c:5}, with other pin numbers, and of
\file{main.c}, with plain (non-\code{static}) functions and no Doxygen comments. The \emph{code
shown in the text} is what you should actually type in; the IDE mechanics (menus, dialogs, build steps) are otherwise
identical.

\subsection{Wire up the External Tool}
\label{c2:c:8}
Programming the chip externally requires the Advanced UI profile: \code{Tools → Select Profile},
choose \textbf{Advanced}, then \textbf{Apply}.

\bookimage[0.95]{lectures/L02/appendix/images/19_select_profile_menu.png}{Tools, Select
Profile.}{c2:fig:c19}

\bookimage[0.95]{lectures/L02/appendix/images/20_advanced_profile_dialog.png}{The Advanced
profile.}{c2:fig:c20}

\noindent Now create the External Tool entry: \code{Tools → External Tools}.

\bookimage[0.95]{lectures/L02/appendix/images/21_external_tools_menu.png}{Tools, External
Tools.}{c2:fig:c21}

\noindent In the dialog that appears, click \textbf{Add}. Title it \code{Arduino} followed by the
COM port used earlier in the Arduino IDE: e.g.\ \code{Arduino - COM4}. Paste the first line from
\code{Arduino via Microchip Studio.txt} into \textbf{Command}, and the second line into
\textbf{Arguments}. Check \textbf{Use Output Window}, then click \textbf{OK}.

\bookimage[0.95]{lectures/L02/appendix/images/22_external_tools_arduino_com4.png}{The External
Tools dialog with the Arduino - COM4 entry: its Command and its Arguments.}{c2:fig:c22}

\noindent Go back into \code{Tools} and select the tool you just created (e.g.\
\code{Arduino - COM4}).

\bookimage[0.95]{lectures/L02/appendix/images/23_tools_menu_arduino_com4_selected.png}{Running the
new tool from the Tools menu.}{c2:fig:c23}

\subsection{Flash and verify}
\label{c2:c:9}
If everything worked, the Output window should end with the number of bytes written, e.g.:

\begin{console}
Writing | ################################################## | 100% 0.05s

avrdude.exe: 194 bytes of flash written
avrdude.exe: safemode: lfuse reads as 0
avrdude.exe: safemode: hfuse reads as 0
avrdude.exe: safemode: efuse reads as 0
avrdude.exe: safemode: Fuses OK (E:00, H:00, L:00)

avrdude.exe done. Thank you.
\end{console}

\bookimage[0.95]{lectures/L02/appendix/images/24_output_window_flash_success.png}{A successful
flash, as the Output window reports it.}{c2:fig:c24}

\noindent The LED on D9 should now light up while the button on D13 is held, and stay off otherwise.
If the write fails instead, with something like:

\begin{console}
avrdude.exe: ser_open(): can't open device "\\.\COM4": Det går inte att hitta filen.
\end{console}

\noindent (the last part is Swedish Windows for ``The system cannot find the file specified''),
check that the board is connected, then check the Arduino IDE for its current COM port; if it's
changed, see below.

\subsubsection{If the COM port changes}
Reconnecting the board sometimes assigns a different COM port, which surfaces as the same
\code{ser_open()} error above. The fix is a second External Tool entry, identical to the first
except for the port: back in \code{Tools → External Tools}, add a new entry (e.g.\
\code{Arduino - COM3}), reusing the same Command, and an Arguments value with \code{PCOM4} swapped
for the new port (e.g.\ \code{PCOM3}).

\bookimage[0.95]{lectures/L02/appendix/images/25_external_tools_menu_for_com3.png}{Back to Tools,
External Tools, for a second entry.}{c2:fig:c25}

\bookimage[0.95]{lectures/L02/appendix/images/26_external_tools_arduino_com3.png}{The Arduino - COM3
entry: the same Command, with \code{-PCOM3} in its Arguments.}{c2:fig:c26}

\noindent Then select the new tool from \code{Tools} and run it the same way as before.

\bookimage[0.95]{lectures/L02/appendix/images/27_tools_menu_arduino_com3_selected.png}{Running the
COM3 tool from the Tools menu.}{c2:fig:c27}
